\scanpage{93}
amené à supposer que \deleted{\unsure{le|les}} $p$ est \unsure{égal\textsuperscript{t}} cofibrant,
i.e. que les adjoints \uline{à gauche} des $f^*$ existent,
[qu'il n'y a pas lieu alors de noter $f_*$, mais
plutôt $f_!$], \inserted{si les $f^*$ sont associés à des \uline{morphismes} $\unsure{f_{!}|f:}$,}
On sait qu'il revient au $=$ de
dire que les $f^*$ commutent aux \inserted{petites} $\underleftarrow{\lim}$,
-- -- ce qui implique alors que ces dernières
se calculent fibre par fibre.

\deleted{\unsure{contre exemple} pour le cas où les $f^*$ ne commutent
pas aux $\underleftarrow{\lim}$ quelconques ?}

Dans le cas général, \unsure{où} $B$ est \deleted{\unsure{équivalente}} ``ess petite'',
il n'est cependant pas difficile de prouver
qu'il y a un petit ens de générateurs
-- à rédiger en détail ..... . Donc on
trouve bien des $\mathcal U$-topos resp $\mathcal U$-ab-topos.
Il y a donc des $\underleftarrow{\lim}$, mais leur
calcul en général, ne se faisant pas
terme à terme, est un peu mystérieux.

% archive sheet number: 92
